\section{结论}
\label{sec:conclusion}

\ours 把精度花在接触所依赖之处：它精确建模目标的度量运动、粗略建模整个场景，并只用这些预测来训练它据以行动的表示。在 \bench、DOMINO 与六个真机任务上，它都优于反应式、动态与世界—动作策略。\fillin{消融表明，目标历史在附加延迟下作用最大，未来监督在不改变推理计算的情况下改善控制。}\fillin{说明专用 SE(3)、场景流与统一精度对照支持什么：更好的预测设计，或者在任务分组也一致时，按角色分配精度。}计入推理延迟后，\fillin{说明成功率与相对表现与暂停评测相比有何不同}。

\paragraph{局限性。}
\ours 需要一个指定的刚体目标及其可靠的位姿估计，其预测也不建模机器人自身的接触如何改变目标的运动。目标轨迹预测头在每个时域只预测一个未来，不表示不确定性或多模态运动；不确定性感知的预测可以弥补这一点。
